\chapter{Explainer: algebra and coalgebra}

\section{Algebra}

% Before we look at coalgebras, let us consider algebras, which are slightly more intuitive.
In algebra, there are many interesting structures to study, such as groups, rings, vector spaces, etc.
If we look at the definitions for some common algebraic structures:

\begin{itemize}
    \item A group is a tuple $(G, \cdot, (-)^{-1})$ satisfying some axioms.
    \item A ring is a tuple $(R, +, \cdot)$ satisfying some axioms.
    \item A boolean algebra is a tuple $(B, \land, \lor, \neg)$ satisfying some axioms.
\end{itemize}

They all look similar -- a carrier, usually a set, followed by some operations on that carrier, satisfying some axioms.
It may be useful to look for an abstraction over all of these.

In order to make it easier to think about abstractions for a universal algebra, we'll start by focusing our work in the category of sets $\Set$.
So our carrier will be some object in $\Set$.
Then, consider an endofunctor $F : \Set \rightarrow \Set$.
This will have two parts: an action on objects, which is simply a set function, and an action on morphisms.
If we take $X$ to be the input carrier set, then $F(X)$ is essentially the representation of \emph{inputs} to operations you can take with $X$.
Then, we have a morphism $\alpha : F(X) \rightarrow X$, known as the \emph{structure} map, which defines how the operations take place.

For example:
\begin{itemize}
    \item the identity element of a group is not dependent on anything, so we say $F(X) = 1$, and so the morphism $\alpha : 1 \rightarrow X$ simply selects the identity element with no further input.
    \item the inverse operation of a group is a set-function $X \rightarrow X$. This translates to our framework as $F(X) = X$ and so $\alpha : X \rightarrow X$ in this case.
    \item the group operation is a binary set-function $(\cdot) : X \times X \rightarrow X$. We have $F(X) = X \times X$, so the morphism $\alpha : X \times X \rightarrow X$ has the correct type.
\end{itemize}

But each of these are still independent pieces of a group definition.
We must combine them somehow.
A single morphism $\alpha : F(X) \rightarrow X$ must represent all actions you can possibly take on a group.
One answer would be to take $F(X) = 1 \amalg X \amalg X \times X$, where $\amalg$ is the disjoint set-union operation.
This is usually written instead as $F(X) = 1 + X + X \times X$, as we will soon generalize disjoint unions in $\Set$ to coproducts in an arbitrary category.

So our combined morphism $\alpha : F(X) \rightarrow X$ essentially gives us a choice of which operation we would like to take, as well as how each is implemented.
It should be clear at this point that each functor $F$ can possibly have multiple $\alpha : F(X) \rightarrow X$.
This corresponds to the notion that there are multiple groups, multiple rings, multiple Boolean algebras, etc.

Of course, since $F$ is a functor, it must have an action on morphisms as well.
However, this action is mostly forced by the universal property of coproducts.
For example, in the case of $F(X) = A \times X + B$, it must hold that $F(A_1 \times X + B_1 \xrightarrow{f} A_2 \times X + B_2)$ must also be a sum, which means it can be decomposed into two functions $f_1 : A_1 \times X \rightarrow A_2 \times X$, and $f_2 : B_1 \rightarrow B_2$.
We will write this as $F(f) = f_1 + f_2$ or sometimes even $[f_1, f_2]$ for short.

To sum this up, here is the formal definition for $F$-algebra:

\begin{definition}
  Let $\calC$ be an ambient category.
  An $F$-algebra is a pair $(X, \alpha : F(X) \rightarrow X)$, where $X \in Ob(\calC)$, $F$ is an endofunctor $\calC \rightarrow \calC$, and $\alpha : F(X) \rightarrow X$.
\end{definition}
